\chapter*{Lecture07整区}
\addcontentsline{toc}{chapter}{Lecture07整区}
\stepcounter{chapter}

\section{分式域}
我们已经知道，从整数可以构建有理数域。事实上这个过程可以进行推广。\\
\kfj{定义：}对于R是一个整区，记0和1分别是R的零元和幺元。对于
\[
R\times R^*=\{(r,r^*)|r\in R,\; r^*\in R^*\}
\]
我们在$R\times R^*$上定义等价关系：
\[
(a,b)\sim(c,d)\text{当且仅当}ad=bc
\]
我们稍后会验证这确实是一个等价关系，此时令$F=(R\times R^*)/\sim$，并记$(a,b)$的等价类为$\dfrac{a}{b}$。\\
我们在F上定义运算：
\[
\left\{
\begin{aligned}
&(1)(\text{加法}) +\colon F\times F\to F ,\; (\frac{a}{b},\frac{c}{d})\mapsto \frac{ad+bc}{bd}\\
&(2)(\text{乘法}) \cdot\colon F\times F\to F ,\; (\frac{a}{b},\frac{c}{d})\mapsto \frac{ac}{bd}
\end{aligned}\right.
\]
我们稍后会验证上述运算良定义，并且F在上述运算下构成一个域。我们称F是R的分式域。\\
\noteb{(注：我们要搞清楚的是，一个环的分式域只有一个（就是上述过程所构造的域）。很多时候我们会遇到环R可以直接作为一个更大域F的子环，且这个更大的域F正好同构于R的分式域。事实上这个时候F一定不是R的分式域)}\\
\kfj{命题：}\\
\indent (1):上述定义中，“$\sim$”是一个等价关系。\\
\indent (2):上述定义中，F上的两个运算良定义。\\
\indent (3):F在上述运算下构成一个域。\\
\begin{proof}
(1):(自反性):由于R是交换环，于是有$\forall (a,b)\in R\times R^*,\; ab=ba$，即$(a,b)\sim (a,b)$。\\
\noteb{这里我们也可以看出，如果你希望一个环可以从他构做分式域，那么交换性是必要条件。}\\
(2):(对称性):$(a,b)\sim (c,d)\iff ad=bc\iff bc=ad \iff (c,d)\sim (a,b)$\\
(3):(传递性):若$(a,b)\sim (c,d),\; (c,d)\sim (x,y)$，则有：$ad=bc$，$cy=dx$，我们希望证明$ay=bx$。这个时候我们可以做这样一个操作：
\[
ad=bc\implies ady=bcy\implies ady=bdx \implies ayd=bxd\implies (ay-bx)d=0
\]
其中，由于$d\in R^*$，且R是整区，特别地R是整的，因此由消去律：
\[
(ay-bx)d=0\implies ay-bx=0\iff ay=bx\iff (a,b)\sim(x,y)
\]
\noteb{从这里我们可以看出，若想构做分式域，则这个环必须是整的。}\\
于是我们验证了上述定义中，“$\sim$”是一个等价关系。\\
(2)(加法):设$(a,b)\sim (a',b'),\; (c,d)\sim (c',d')$。于是有：$ab'=ba',\; cd'=dc'$，且依定义有：
\[
\frac{a}{b}+\frac{c}{d}=\frac{ad+bc}{bd},\; \frac{a'}{b'}+\frac{c'}{d'}=\frac{a'd'+b'c'}{b'd'}
\]
我们从等价关系的定义去验证：
\[
(ad+bc)b'd'-bd(a'd'+b'c')=adb'd'+bcb'd'-bda'd'-bdb'c'=(adb'd'-bda'd')+(bcb'd'-bdb'c')=0
\]
因此可知:
\[
\frac{ad+bc}{bd}=\frac{a'd'+b'c'}{b'd'}
\]
即上述定义的加法是良定义的。\\
(乘法):设$(a,b)\sim (a',b'),\; (c,d)\sim (c',d')$。于是有：$ab'=ba',\; cd'=dc'$，且依定义有：
\[
\frac{a}{b}\cdot \frac{c}{d}=\frac{ac}{bd},\; \frac{a'}{b'}\cdot \frac{c'}{d'}=\frac{a'c'}{b'd'}
\]
我们仍从等价关系的定义去验证：
\[
acb'd'-bda'c'=0\implies \frac{ac}{bd}=\frac{a'c'}{b'd'}
\]
即上述定义的乘法是良定义的。\\
(3)此时我们只需证明$(F,+)$和$(F^*,\cdot)$都构成Abel群即可。\\
$(F,+)$:由于R是整区，于是有：
\[
\forall \frac{a}{b},\frac{c}{d},\frac{x}{y}\in F,\; (\frac{a}{b}+\frac{c}{d})+\frac{x}{y}=\frac{(ad+bc)y+bdx}{bdy}
\]
而
\[
\frac{a}{b}+(\frac{c}{d}+\frac{x}{y})=\frac{ady+b(cy+dx)}{bdy}=\frac{(ad+bc)y+bdx}{bdy}
\]
因此可知加法具有结合律。易知：
\[
\forall \frac{a}{b}\in F,\; \frac{a}{b}+\frac{0}{1}=\frac{a}{b}=\frac{0}{1}+\frac{a}{b},\; \frac{a}{b}+\frac{-a}{b}=0
\]
因此F在加法下构成一个群，并且再由R的交换性可知：
\[
\forall \frac{a}{b},\frac{c}{d}\in F,\; \frac{a}{b}+\frac{c}{d}=\frac{ad+bc}{bd}=\frac{cb+da}{db}=\frac{c}{d}+\frac{a}{b}
\]
因此$(F,+)$是一个Abel群。\\
$(F^*,\cdot)$:我们由R中乘法的结合律立即得到：
\[
\forall \frac{a}{b},\frac{c}{d},\frac{x}{y}\in F,\;(\frac{a}{b}\cdot\frac{c}{d})\cdot\frac{x}{y}=\frac{(ac)x}{(bd)y}=\frac{a(cx)}{b(dy)}=\frac{a}{b}\cdot (\frac{c}{d}\cdot \frac{x}{y})
\]
故$F^*$中乘法具有结合律。易知：
\[
\forall \frac{a}{b}\in F^*,\; \frac{a}{b}\cdot \frac{1}{1}=\frac{a}{b}=\frac{1}{1}\cdot\frac{a}{b},\; \frac{a}{b}\cdot \frac{b}{a}=\frac{1}{1}=\frac{b}{a}\cdot\frac{a}{b}
\]
因此$F^*$在此乘法下构成一个群。再由R的交换性：
\[
\forall \frac{a}{b},\frac{c}{d}\in F^*,\; \frac{a}{b}\cdot \frac{c}{d}=\frac{ac}{bd}=\frac{ca}{db}=\frac{c}{d}\cdot\frac{a}{b}
\]
于是$(F^*,\cdot)$是一个Abel群。综上可知，F是一个域。
\end{proof}
\kfj{记号：}在上述证明中，我们知道$\dfrac{0}{1}$和$\dfrac{1}{1}$分别是F的零元和幺元，那么我们约定在分式域中，就记$\dfrac{0}{1}$为0，记$\dfrac{1}{1}$为1。\\
\kfj{定理：}对于整区R，设F是R的分式域，映射$i\colon R\to F,\; r\mapsto \frac{r}{1}$称为环R到其分式域的标准嵌入。任取一个环同态$g\colon R\to F'$，且满足$F'$是域，则使得下图交换的环同态$h\colon F\to F'$存在且唯一。
\[
\begin{tikzcd}[row sep=large, column sep=large]
R \arrow[r, "g"] \arrow[d, "i"'] & F' \\
F \arrow[ru, dashed, "h"'] &
\end{tikzcd}
\]
\begin{proof}
(存在性):做$h\colon F\to F',\frac{a}{b}\mapsto g(a)\inv{[g(b)]}$。我们先证明$h$是良定义的。\\
设$(a,b)\sim (a',b')$，则有$ab'=ba'$\noteb{(注：R中元素不可直接取逆)}。由于$F'$是域，故有：\\
\[
g(ab')=g(ba')\iff g(a)g(b')=g(b)g(a')\iff g(a)\inv{[g(b)]}=g(a')\inv{[g(b')]}\iff h(\frac{a}{b})=h(\frac{a'}{b'})
\]
于是可知$h$是良定义的映射。\\
(再证$h$是群同态)设$\dfrac{a}{b},\dfrac{c}{d}\in F$，则有：
\[
h(\dfrac{a}{b}+\dfrac{c}{d})=h(\frac{ad+bc}{bd})=g(ad+bc)\inv{[g(bd)]}
\]
由于$g\colon R\to F'$是环同态，因此有：
\[
h(\dfrac{a}{b}+\dfrac{c}{d})=[g(a)g(d)+g(b)g(c)]\inv{[g(b)]}\inv{[g(d)]}=g(a)\inv{[g(b)]}+g(c)\inv{[g(d)]}=h(\frac{a}{b})+h(\frac{c}{d})
\]
因此可知$h$是加群同态。\\
(再证$h$是环同态):设$\dfrac{a}{b},\dfrac{c}{d}\in F$，由于$F'$是域，则有：
\[
h(\frac{a}{b}\cdot \frac{c}{d})=h(\frac{ac}{bd})=g(ac)\inv{[g(bd)]}=g(a)\inv{[g(b)]}\cdot g(c)\inv{[g(d)]}=h(\frac{a}{b})\cdot h(\frac{c}{d})
\]
因此可知$h$是环同态，此时易知$h$满足$hi=g$，即所写交换图交换。\\
(唯一性):设$f\colon F\to F'$是环同态，且满足$fi=g$。\\
那么可知：$\forall \dfrac{r}{1}\in F,\; f(\dfrac{r}{1})=fi(r)=g(r)$\\
而由于$f$是环同态可知，$\forall \dfrac{r}{1}\in F^*,\; f(\dfrac{1}{r})=f(\inv{\left(\dfrac{r}{1}\right)})=\inv{f(\dfrac{1}{r})}=\inv{[g(r)]}$。\\
于是我们有：
\[
\forall \frac{a}{b}\in F,\; f(\frac{a}{b})=f(\frac{a}{1}\cdot \frac{1}{b})=f(\frac{a}{1})\cdot f(\frac{1}{b})=g(a)\inv{[g(b)]}
\]
于是可知，满足该条件的环同态唯一。
\end{proof}
\indent 上述定理有时也称为“分式域的泛性质”，泛性质一般都包含“存在性”和“唯一性”两个关键的性质，它们会帮我们简洁又漂亮地解决一些问题。下面我们会见到这样的“泛性质”将会有怎样强大的功能。\\
\kfj{例：}$\mathbb{Z}$的分式域与$\mathbb{Q}$同构。
\begin{proof}
设$F$是$\mathbb{Z}$的分式域，我们直接构造映射：
\[
\gamma \colon \mathbb{Q}\to F,\; \dfrac{p}{q}\mapsto \dfrac{p}{q}
\]
\noteb{(注：上述映射两侧的分式具有不同的含义，$F$中的分式意为等价类，$\mathbb{Q}$中的分式意为有理数。)}\\
\indent 先证明$\gamma$是良定义的。在$\mathbb{Q}$中，有:
\[
\frac{p}{q}=\frac{p'}{q'}\iff pq'=qp'\iff pq'-qp'=0
\]
\indent 其中第一个等价符号是由于$\mathbb{Q}^*$在乘法下构成一个群从而有消去律；第二个等价符号是由于$\mathbb{Q}$具有环结构因而可以取加法逆元。而依定义在$F$中有：
\[
\frac{p}{q}=\frac{p'}{q'}\iff pq'-qp'=0
\]
\indent 因此可知$\gamma$是一个良定义的单射。于是从$\gamma$的形式来看易知其是一个满射，再由$F$和$\mathbb{Q}$上的运算可知$\gamma$是一个群同态，从而是一个群同构。
\end{proof}

\kfj{命题：}设$R$是域，则依定义知$R$一定是整区。于是设$F$是$R$的分式域，则有结论：
\[
\phi\colon F\to R,\text{是环同构}
\]
\begin{proof}
\indent 对于标准嵌入$I\colon R\to F$，我们应用分式域的泛性质，并直接取定理中的“$F'$”为$F$本身。于是可知使得下图交换的环同态$h\colon F\to F'$存在且唯一：
\[
\begin{tikzcd}[row sep=large, column sep=large]
R \arrow[r, "i"] \arrow[d, "i"'] & F \\
F \arrow[ru, dashed, "h"'] &
\end{tikzcd}
\]
\indent 而事实上，恒等映射$\mathrm{Id}_F\colon F\to F$是环同态且满足：$\mathrm{Id}_F\circ i=i$，因此由唯一性知：必有$h=\mathrm{Id}_F$。这其实告诉我们此时对于映射$i$的复合具有消去律，即：若环同态$h\colon F\to F$，满足$hi=i$，则必有$h=\mathrm{Id}_F$。\\
\indent 下面我们再取$F'$就是$R$，这是由于此时$R$是一个域。以及环同态$\mathrm{Id}_R\colon R\to R $是R上的恒等同态。于是再由泛性质可得：使得下图交换的环同态$\phi\colon F\to R$存在且唯一。
\[
\begin{tikzcd}[row sep =large, column sep =large]
 & F\arrow[d,dashed,"\phi"']\\
R\arrow[ru,"i"]\arrow[r,"\mathrm{Id}_R"] & R
\end{tikzcd}
\]
\indent 于是我们可知：$\phi i=\mathrm{Id}_R$。易知恒等映射满足：$i\circ \mathrm{Id}_R=i$，因此也有交换图：
\[
\begin{tikzcd}[row sep = large, column sep = large]
R\arrow[r,"\mathrm{Id}_R"]\arrow[rd,"i"'] & R\arrow[d,"i"]\\
 & F
\end{tikzcd}
\]
我们把上述交换图拼在一起得到：
\[
\begin{tikzcd}[row sep =large, column sep =large]
 & F\arrow[d,dashed,"\phi"]\\
R\arrow[ru,"i"]\arrow[r,"\mathrm{Id}_R"]\arrow[rd,"i"'] & R\arrow[d,"i"]\\
 & R
\end{tikzcd}
\]
\indent 这个时候我们沿着这个图的最外圈，也就是左边的$R$通过$i$指向F，再通过$\phi$和$i$向下指向$R$，这条路径应该等价于从左边的$R$直接指向右下角的$R$。于是可知：$i\phi\circ i=i\implies i\phi=\mathrm{Id}_F$ 。我们综合上面得到的关系：
\[
\left\{
\begin{aligned}
&\phi i=\mathrm{Id}_R\\
&i \phi=\mathrm{Id}_F
\end{aligned}\right. 
\]
因此可知环同态$\phi\colon F\to R$是环同态。
\end{proof}
\kfj{推论：}互不同构的整区可以有互相同构的分式域。
\begin{proof}
可知$\mathbb{Z}$不是域，$\mathbb{Z}$的分式域与$\mathbb{Q}$同构。而由命题知$\mathbb{Q}$的分式域与其自身同构。因此此时有$\mathbb{Z}$与$\mathbb{Q}$不同构，但是他们的分式域是互相同构的。
\end{proof}

\section{唯一析因环UFD}

\kfj{定义：}设R是整区，则给出如下定义：\\
\indent (1)：R中所有单位(unit)，即可逆元，所组成的集合记为$U$。易知$U$在环$R$的乘法下构成一个Abel群，我们称之为环$R$的单位群。\\
\indent (2)：对于$R^*$中元素$r$，若存在$a,b\in R^*$，使得$a\cdot b=r$，则称$a$是$r$的因子，称$a$整除$r$，记为$a|r$。（同理$b$也是$r$的因子，也有$b|r$）\\
\indent (3)：$\forall r\in R^*$，对于单位群中的元素$u$，可知：$r=(u\cdot \inv{u})r=u(\inv{u}\cdot r)\implies u|r$。这件事对于所有的元素$r$都成立，因此此时$u$称为$r$的一个平凡因子。\\
\indent (4)：对于$a,b\in R^*$，若$a|b$且$b|a$，则称$a$和$b$是相伴的，记为$a\sim b$。\\
\kfj{性质：}
\[
\left\{
\begin{aligned}
&\text{1、}\forall a,b\in R^*,\; a\sim b\iff \exists u\in U,\text{使得}ua=b\\
&\text{2、"}\sim\text{"是一个等价关系，且}a\sim b,\; c\sim d\text{同时成立意味着}ac\sim bd\\
&\text{3、}u\in U\iff u\sim 1
\end{aligned}\right.
\]
\begin{proof}
\indent 1、($\impliedby$)此时有$ua=b$和$a=\inv{u}b$，因此$a|b,\; b|a$，因而有$a\sim b$\\
($\implies$)此时可知$a|b$且$b|a$，即$\exists x,y\in R^*$使得$xb=a,\; ya=b$。于是有$y(xb)=b$，即$(yx)b=b$。在整区中有消去律，于是得到$yx=1$。由交换性有$xy=yx=1$，因此$x$和$y$都是可逆元。由此我们得到：$ya=b,\; y\in U$。\\
\indent 2、由性质1立即可知$"\sim"$是一个等价关系，下面我们来看对于$a\sim b$且$c\sim d$，我们设$u,v\in U$，使得$ua=b,\; vc=d$。于是由整区的交换性，$bd=(ua)(vc)=u(av)c=u(va)c=(uv)(ac)\implies bd\sim ac$\\
\indent 3、由性质1可知$u\sim 1\iff \exists v\in U$使得$uv=1$。而$u\in U\iff \exists v\in R^*$使得$uv=1$。事实上若$uv=1$，则$u$和$v$都可逆，因此必有$u,v\in U$，于是可知$\exists v\in R^*$使得$uv=1$$\iff \exists v\in U$使得$uv=1$。综上可知有$u\sim 1\iff u\in U$
\end{proof}
\kfj{性质：}设$R$是整区，$U$是$R$的单位群，那么：
\[
\left\{
\begin{aligned}
&(1):\;a\mid b\iff <b>\subset <a>\\
&(2):\;a\sim b\iff <a>=<b>\\
&(3):\;u\in U\iff <u>=R
\end{aligned}\right.
\]
\begin{proof}
(1):可知$a\mid b\iff \forall r\in R^*,\; a\mid rb$。（其中左推右显然，右推左令$r$是幺元即可）然而易知$0\in<b>\cap<a>$，因此$\forall r\in R^*,\; a\mid rb\iff <b>\subset<a>$，从而有：$a\mid b\iff <b>\subset<a>$\\
\indent (2):可知$a\sim b\iff a\mid b\text{且} b\mid a$，因此由性质(1)可知：
$\left\{
\begin{aligned}
&<b>\subset<a>\\
&<a>\subset<b>
\end{aligned}\right.$
\\
于是立即可知$a\sim b\iff <a>=<b>$。\\
\indent (3):设环$R$的幺元是$1$，那么有：
\[
<1>=\{r\cdot 1|r\in R\}=R
\]
因此借助(2)的结论可知$<u>=R\iff <u>=<1>\iff u\sim 1\iff u\in U$
\end{proof}
\kfj{定义：}
\[
\begin{aligned}
&(1)\text{对于}a,b\in R^*,\; \text{若}a|b,\text{但是}b\nmid a,\; \text{则称}a\text{是}b\text{的真因子}\\
&(2)\text{对于}R^*-U\text{中元素$r$，若$r$的真因子只有单位(unit)，则称$r$是不可约元素，否则称$r$是可约元。}\\
&(3)\text{对于}R^*-U\text{中元素$p$，若对任意}a,b\in R^*\text{满足}p|a\cdot b\text{都可以得出}p|a\text{或}p|b,\; \text{则称p是素元素}\\
\end{aligned}
\]
\noteb{(注：其中素元素的定义也可以陈述为：$p\nmid a$且$p\nmid b\implies p\nmid a\cdot b$)}\\
\kfj{命题：}a的真因子是与a不相伴的因子（易证）\\
\kfj{命题：}素元素是不可约元素
\begin{proof}
\indent 设$p$是素元素，且$a$是$p$的真因子，则依定义有$a\mid p$，但$p\nmid a$。此时我们设$p=ra$，其中$r\in R^*$。于是有$p\mid ra$，再由于$p$是素元素，且$p\nmid a$，于是有$p\mid r$。此时再由定义可知：$\exists v\in R^*$，使得$vp=r$，从而有$vpa=p\iff p(va-1)=0$。由于整区中有消去律，因此可知$va=1$，即$a\in U$。因此可知$p$是不可约元素。
\end{proof}
\kfj{命题：}不可约元素不一定是素元素
\begin{proof}
考虑环：
\[
\mathbb{Z}[\sqrt{-5}]=\{a+b\sqrt{-5}|a,b\in \mathbb{Z}\}
\]
(Step1):先证明$\mathbb{Z}[\sqrt{-5}]$是整区\\
考虑其中减法：
\[
(a+b\sqrt{-5})-(c+d\sqrt{-5})=(a-c)+(b-d)\sqrt{-5}\in \mathbb{Z}[\sqrt{-5}]
\]
考虑其中乘法：
\[
(a+b\sqrt{-5})\cdot (c+d\sqrt{-5})=(ac-5bd)+(bc+ad)\sqrt{-5}\in \mathbb{Z}[\sqrt{-5}]
\]
于是可知$\mathbb{Z}[\sqrt{-5}]$在复数乘法下构成一个环。同时有：
\[
(a+b\sqrt{-5})\cdot (c+d\sqrt{-5})=0\iff (ac-5bd)+(bc+ad)\sqrt{-5}=0
\]
\indent 于是$ac=5bd,\; ad=-bc$。前式同时乘$d$得：$acd=5bd^2\iff 5bd^2+bc^2=0\iff b(5d^2+c^2)=0$。于是由整数的性质知，此时有要么$b=0$，要么$c=d=0$。\\
\indent 若$c$和$d$不同时为0，则有$b=0$，于是带入原式有：
\[
a(c+d\sqrt{-5})=0\iff ac+ad\sqrt{-5}=0\iff ac=ad=0\implies a=0
\]
此时则可知$a=b=0$，即$a+b\sqrt{-5}=0$。\\
\indent 若$c=d=0$，则$c+d\sqrt{-5}=0$。于是可知$\mathbb{Z}[\sqrt{-5}]$是整的。易知$\mathbb{Z}[\sqrt{-5}]$是交换环，且幺元为整数$1$，因此综上可知$\mathbb{Z}[\sqrt{-5}]$是整区。\\
(Step2):求$\mathbb{Z}[\sqrt{-5}]$的单位群$U$\\
\indent 取$\alpha\in U$，不妨设$\alpha=x+y\sqrt{-5}$，我们取复数模长有：${|\alpha|}^2=x^2+5y^2$。由于$\alpha\in U$，故存在$\beta\in {\mathbb{Z}[\sqrt{-5}]}^*$使得$\alpha\cdot \beta=1$，于是由复数取模长性质得：
\[
{|\alpha\cdot \beta|}^2=1\iff {|\alpha|}^2\cdot{|\beta|}^2=1
\]
\indent 此时由于${|\alpha|}^2$和${|\beta|}^2$都是正整数，因此有${|\alpha|}^2={|\beta|}^2=1$，特别地，有$x^2+5y^2=1$。事实上，若$y\neq 0$，则可知$x^2+5y^2=\geq 5$，不合题。因此必有$y=0$，从而有$a=\pm 1$，即$\alpha=\pm 1,\; \implies U\subset\{1,-1\}$。而不难证明$1$和$-1$都是$\mathbb{Z}[\sqrt{-5}]$中得可逆元，因此有：$U=\{1,-1\}$。\\
(Step3):证明整数$3$在$\mathbb{Z}[\sqrt{-5}]$中是不可约元素，但不是素元素。\\
\indent 设$\alpha$是$3$的一个真因子，且有$\alpha\cdot \beta=3$，设$\alpha=a+b\sqrt{-5}$。则可知：
\[
{|\alpha|}^2\cdot{|\beta|}^2={|\alpha\cdot \beta|}^2=3^2=9
\]
于是由整数性质知：${|\alpha|}^2=1$或$3$或$9$。\\
(${|\alpha|}^2=9$时)，有整数${|\beta|}^2=1$，因此$\beta\in U\implies \alpha\sim 3$，故$\alpha$不是$3$的真因子。\\
(${|\alpha|}^2=3$时)，有$a^2+5b^2=3$。事实上当$b=0$时，有$a^2=3$在整数中无解（$|a|\geq 2$时 $a^2\geq 4>3$；$0\leq |a|\leq 1$时 $a^2\leq 1<3$）；当$b\neq 0$时，$a^2+5b^2\geq 5>3$因此也无解，故此时不合题。\\
(${|\alpha|}^2=1$时)，有$\alpha\in U$，是单位。\\
\indent 于是我们知道，$3$的真银子只能是单位，故$3$是不可约元。而观察可知：
\[
3\times 3=9 =(2+\sqrt{-5})(2-\sqrt{-5})\implies 3|(2+\sqrt{-5})(2-\sqrt{-5})
\]
\indent 而若$3|(2+\sqrt{-5})$，则$\exists \beta\in {\mathbb{Z}[\sqrt{-5}]}^*$，使得$3\cdot \beta=(2+\sqrt{-5})$。等式两端同时取模长得：
\[
{|3\beta|}^2={|2+\sqrt{-5}|}^2\iff 9{|\beta|}^2=9\iff {|\beta|}^2=1
\]
\indent 于是有$\beta\in \{1,-1\}$，因而$3=2+\sqrt{-5}$或$3=-2-\sqrt{-5}$。变形得$1=\sqrt{-5}$或$5+\sqrt{-5}=0$，可知两种情况都矛盾，因此$3\nmid (2+\sqrt{-5})$\\
\indent 同理有$3\nmid (2-\sqrt{-5})$，因此可知$3$不是素元素\\
\end{proof}

\kfj{定义：}对于整区$R$，我们作如下定义：\\
\indent (素性条件)：若$R$中不可约元素都是素元素，则称$R$满足素性条件。\\
\indent (最大公因子条件)：对于$R^*$中元素$a$和$b$，若$\exists d\in R^*$，使得$d|a$且$d|b$，则称$d$是$a$和$b$的公因子。进一步地，若对于$a$和$b$的任意一个公因子$d'$，都有$d'|d$，则称$d$是$a$和$b$的一个最大公因子。于是，若$R^*$中任意两个元素$a$和$b$，都存在他们的最大公因子，则称整区$R$满足最大公因子条件\\
\indent \noteb{(注：与最大公因子相伴的元素也是最大公因子，因此不能确定最大公因子唯一)}\\
\indent (因子链条件)：对于$R^*$中的一个无穷序列：
\[
a_1,\; a_2,\; \cdots ,\;a_n,\;a_{n+1},\cdots
\]
若$\forall n\geq 1$，都有$a_{n+1}|a_n$，则称这个序列是一个因子链。若整区$R$满足对于$R$中的任意一个因子链${\{a_n\}}_{n\in \mathbb{N}^*}$都有：
\[
\exists N\in \mathbb{N}^*,\; \forall n\geq N,\; a_n\sim a_N
\]
则称$R$满足因子链条件。\\
\indent (有限析因条件)：若$R-U$中任意一个元素都可以写为有限个不可越元的乘积，那么则称$R$满足有限析因条件\\
\kfj{定义：}(唯一析因环)：对于整区$R$，若$R$满足有限析因条件，并且任取一个元素$r\in R^*$，若$R$中有限个不可约元素$p_1,\; p_2,\cdots,p_s$，满足$p_1p_2\cdots p_s=r$，则称这是一个$r$的有限不可约乘积表示。若任取两个有限不可约乘积表示：
\[
r=p_1p_2\cdots p_s,\; r=q_1q_2\cdots q_t
\]
都满足$s=t$，且存在置换$\sigma\in S_t$，使得$\forall 1\leq i\leq t,\; p_i\sim q_{\sigma(i)}$。则称$R$是一个唯一析因环，简称UFD(Unique Factorization Domain)\\
\kfj{引理：}对与$R$是整区，有：
\[
(\text{因子链条件})\implies (\text{有限析因条件})
\]
\begin{proof}
(先证明$R^*-U$中任意元素都存在不可约的因子)\\
\indent 我们观察因子链条件的陈述不难想到利用它的方法，即反证法然后推出矛盾的无穷序列。这种方法与微积分中很多极限的命题颇有一曲同构之妙，我们下面就来把这个想法具体实施一下。\\
\indent1、任取$r\in R^*-U$，我们假设$r$的因子都是可约元。\\
\indent2、于是由可约元的定义可知，我们可约找到$r$的真因子$r_1\in R^*-U$。我们可知$r_1$的因子也是$r$的因子，因此同理可以找到$r_1$的真因子$r_2$。\\
\indent3、以此类推，我们可以找到一串序列$r,\; r_1,\; r_2,\; \cdots$，且由于$\forall n\in \mathbb{N}^*,\; r_{n+1}|r_n$，因此这是一条因子链，并且$\forall n$，$m>n$时都有$r_m$是$r_n$的真因子，这与因子链条件矛盾。\\
\indent 于是我们可知，$r$存在不可约的因子。\\
(再证明此时满足有限析因条件)\\
\indent1、任取$a_0\in R^*-U$，由刚才论述知$a_0$有不可约因子，于是设$b_1$是$a_0$的一个不可约因子。依定义可知$\exists a_1\in R^*$，使得$a_0=b_1\cdot a_1$。\\
\indent2、若此时$a_1$不可约则直接证毕，于是设$a_1$可约，从而存在$a_1$的不可约因子$b_2$，以及$a_2\in R^*$使得$a_1=b_2\cdot a_2$。此时有$a_0=b_1\cdot b_2\cdot a_2$。\\
\indent 3、此过程可以一直重复。若在有限步之内无法停止，即每次得到的“$a_n$”都是可约的，那么我们得到了一个无穷序列：
\[
a0,\; a_1,\; a_2,\; \cdots
\]
且$\forall n\in \mathbb{N}^*,\; a_n|a_{n+1}$。于是这与因子链条件矛盾，因此事实上这样的操作可以在有限步之内停止，即必然存在$m\in \mathbb{N}^*$，使得上述操作所得$a_m$是不可约元素。于是有有限乘积表示：
\[
a_0=b_1b_2\cdots b_ma_m
\]
\indent 综上可知，此时整区$R$满足有限析因条件
\end{proof}
\indent 我们即将介绍本节最主要的一个定理，但是在证明这个定理之前，我们给出一个概念的定义：\\
\kfj{定义：}设$R$是UFD，我们现在在UFD中定义元素的长度。任取$R^*$中元素$r$，则对于$r=p_1\cdots p_t$是$r$的一个有限不可约乘积表示，则依定义，该乘积表示中的不可约元素的个数是确定的。由此我们定义$r$的长度为$t$，记为$|r|=t$。对于$R$中单位，我们人为定义其长度为$0$；对于$R$中零元我们不定义其长度。\\
\kfj{性质：}
\[
\left\{
\begin{aligned}
&1^*:\;\forall a,b\in R^*,\; |a\cdot b|=|a|+|b|\\
&2^*:\;\forall a,b\in R^*,\; a|b\implies |b|\geq|a|\\
&3^*:\;\forall a,b\in R^*,\; a\sim b\iff (|a|=|b|\text{且}\; a|b)
\end{aligned}\right.
\]
\begin{proof}
$1^*$：由定义不难验证，对于两个有限不可约乘积表示$a=p_1\cdots p_t,\; b=q_1\cdots q_s$，则$ab=p_1\cdots p_tq_1\cdots q_s$是一个$ab$的有限不可约乘积表示，于是有$|ab|=s+t=|a|+|b|$。\\
\indent $2^*$：依定义，$a|b\implies \exists q\in R^*$使得$aq=b$，于是由性质1可知：
\[
|q|\geq 0\implies|b|=|a|+|q|\geq|a|
\]
\indent $3^*$：$a\sim b$推出$|a|=|b|$和$a|b$是显然的，下面证明另一侧的推出关系。已知$a|b$，因此$\exists u\in R^*$，使得$b=ua$。若$u\in R^*-U$，则可知$|u|=1$，因此有：
\[
|b|=|a|+|u|>|a|
\]
与$|a|=|b|$矛盾。因此$u\in U$，即$a\sim b$。
\end{proof}
\kfj{定理：}若$R$是整区，那么下列陈述等价：
\[
\left\{
\begin{aligned}
&(1)R\text{是唯一析因环}\\
&(2)R\text{满足(因子链条件)}+\text{(最大公因子条件)}\\
&(3)R\text{满足(因子链条件)}+\text{(素性条件)}
\end{aligned}\right.
\]
我们分三个部分证明这个定理，即$(1)\implies(2)$、$(2)\implies(3)$、$(3)\implies(1)$三个部分。\\
\kfj{(1)$\implies$(2)}
\begin{proof}
\indent 对于唯一析因环$R$，我们下面先证明其符合因子链条件。设
\[
a_1,\; a_2,\; a_3,\cdots
\]
是$R$中的一条因子链，则可知：
\[
\forall n\in \mathbb{N}^* ,\; |a_n|\geq |a_{n+1}|
\]
若不存在$m\in \mathbb{N}^*$，使得$\forall n\geq m,\; |a_n|=|a_m|$，则有：
\[
\forall n\in \mathbb{N}^*,\; \exists k_n\in\mathbb{N}^*,\; |a_{k_n}|<|a_n|
\]
于是对于UFD中因子链的首项$a_1$应用有限析因条件可知存在正整数$M$，使得$|a_1|<M$。于是可知：
\[
\exists k_1,\cdots,k_M\in \mathbb{N}^*,\; M>|a_1|>|a_{k_1}|>\cdots>|a_{k_M}|
\]
此时有$|a_{k_M}|\leq|a_1|-M<M-M=0$，这与UFD中元素长度的定义相矛盾，因此有：
\[
\exists m\in \mathbb{N}^*,\; \forall n\geq m,\; |a_n|=|a_m|
\]
于是由UFD中长度的性质可知$R$满足因子链条件。\\
\indent 下面证明$R$也满足最大公因子条件。$\forall a,b\in R^*$，若$a$和$b$中有一个是单位，则由于单位长度是0，故立即可知$a$和$b$的公因子必然是单位。\\
若$a$和$b$均不是单位，则有唯一析因：
\[
a=a_1\cdots a_t,\; b=b_1\cdots b_s
\]
将$a_1,\cdots,a_t,b_1,\cdots,b_s$中所有互不相伴的元素取出并记为$p_1,p_2,\cdots,p_k$。即：
\[
\forall 1\leq i<j\leq k,\; p_i\nsim p_j
\]
\indent 且有任取$1\leq i\leq t$，存在唯一的$p_{k_i}$使得$a_i\sim p_{k_i}$；任取$1\leq j\leq t$，存在唯一的$p_{s_j}$使得$b_j\sim p_{s_j}$。于是可知，存在唯一一组系数$n_1,\cdots,n_k,m_1,\cdots,m_k$，使得：
\[
\left\{
\begin{aligned}
&a\sim p_1^{n_1}\cdots p_k^{n_k}\\
&b\sim p_1^{m_1}\cdots p_k^{m_k}
\end{aligned}\right.
\]
\indent 于是对于$\forall 1\leq i\leq k$，令$e_i=min\{m_i,n_i\}$，设$d=p_1^{e_1}\cdots p_k^{e_k}$。此时可知$d$必然是$a$和$b$的公因子。若$d'$也是$a$和$b$的公因子，则$\exists u,b\in R^*$，使得$a=ud',\; b=vd'$。对$\forall 1\leq i\leq k$，若$p_i^{e_i+1}|d'$，则由唯一析因性知，$a$和$b$必有$p_i^{e_i+1}$这个因子。但是$e_i$作为$m_i$和$n_i$中的最小者，必有一个比$e_i+1$小，因此可知$a=ud'$和$b=vd'$中必有一个是矛盾式。因此$\forall 1\leq i\leq k,\; p_i^{e_i+1}\nmid d'$。\\
\indent 再由唯一析因性可知$d'$必然由$\{p_i\}$的乘积而得，故
\[
d'=p_1^{x_1}\cdots p_k^{x_k}
\]
\indent 但是我们由于$p_i^{e_i+1}\nmid d'$，故$\forall 1\leq i\leq k,\; x_i\leq e_i$，于是有$d'|d$，即$d$是$a$和$b$的最大公因子。\\
综上可知，UFD满足即满足因子链条件，也满足最大公因子条件。
\end{proof}
在证明下一个部分之前，我们先证明一个引理：\\
\kfj{引理：}在满足最大公因子条件的整区$R$中，设元素$a$和$b$的最大公因子为$d$，那么任取元素$c\in R^*$，都有$cd$是$ca$和$cb$的一个最大公因子。
\begin{proof}
此时我们不妨设$x$是$ca$和$cb$的一个最大公因子，那么有：
\[
\left\{
\begin{aligned}
&d|a\\
&d|b
\end{aligned}\right.
\implies
\left\{
\begin{aligned}
&cd|ca\\
&cd|cb
\end{aligned}\right.
\]
得到$cd$是$ca$和$cb$的一个公因子，依最大公因子定义可知$cd|x$，不妨设$x=ucd,\; u\in R^*$。而
\[
\left\{
\begin{aligned}
&c|ca\\
&c|cb
\end{aligned}\right.
\implies c|x
\]
\indent 于是此时设$x=cy$。带入$x=ucd$得：$cy=ucd$。由于整区中有消去律，因此可知$y=ud$，带入得$x=cud=ucd$，推出$x|cd$。再结合$cd|x$知$x\sim cd$，因此$cd$也是$ca$和$cb$的一个最大公因子。
\end{proof}
\kfj{(2)$\implies$(3)}
\begin{proof}
\indent 设整区$R$满足因子链条件以及最大公因子条件，任取$R$中不可约元素$r$，下面证明$r$是素元素。\\
\indent 设$a$和$b$满足$r|ab$，其中$a,b\in R^*$。我们记$d$是$b$和$r$的一个最大公因子。由于$r$不可约，因此要么$r\sim d$，要么$d\sim 1$。\\
\indent 对于$(r\sim d)$的情况，由于$d|b$，因此有$r|b$\\
\indent 对于$(d\sim 1)$的情况，由引理：
\[
a=a\cdot 1\text{是}ar\text{和}ab\text{的一个最大公因子}
\]
且有
\[
\left\{
\begin{aligned}
&r|ab\\
&r|ar
\end{aligned}\right.
\implies
r\text{是}ar\text{和}ab\text{的一个公因子}
\]
于是依最大公因子定义可知：$r|a$\\
\indent 综上可知，若$r|ab$，则要么$r|b$，要么$r|a$，因此$r$是素元素。
\end{proof}

\kfj{(3)$\implies$(1)}
\begin{proof}
我们在本定理之前已经证明了(因子链条件)可以推出(有限析因条件)，因此只需证明此时$R$亦满足(唯一析因条件)即可\\
\indent 任取整区$R$中元素$a$，设$a=p_1\cdots p_r$和$a=q_1\cdots q_s$是$a$的两个有限不可约乘积表示，不妨设$s\geq r$。为了方便书写，我们引入除法记号：
\[
\text{记：}q_1\cdots q_s/q_k=\prod_{i\neq k}q_i
\]
\indent 由于$p_1|p_1\cdots p_r\implies p_1|q_1\cdots q_s$，且由素性条件知$p_1$是不可约元因而是素元素，所以$\exists k_1\in\{1,\cdots,s\}$，使得$p_1|q_{k_1}$。而$q_{k_1}$是不可约元且$p_1$不是单位，因此有相伴关系：
\[
p_1\sim q_{k_1}
\]
\indent 由于$p_2|p_2\cdots p_r$，可知$p_2|(q_1\cdots q_s/q_{k_1})$，同理$\exists k_2\in (\{1,\cdots,s\}-\{k_1\})$，使得$p_2\sim q_{k_2}$。\\
\indent 我们可重复此操作：由于$p_i|p_i\cdots p_r$，可知$p_i|(q_1\cdots q_s/q_{k_1}\cdots q_{k_{i-1}})$。同理$\exists k_i\in (\{1,\cdots,s\}-\{k_1,\cdots,k_{i-1}\})$，使得$p_i\sim q_{k_i}$。\\
\indent 此时我们定义映射：
\[
\sigma \colon \{1,\cdots,r\}\to \{1,\cdots,s\},\; i\mapsto k_i
\]
\indent 依定义可知$\forall i\in \{2,\cdots,r\},\; \forall j<i,\; \sigma(i)\neq\sigma(j)$，于是立即可知$\sigma$是单射。同时有：
\[
\{1,\cdots,s\}-\{k_1,\cdots,k_{r}\}=\{1,\cdots,s\}-\sigma(\{1,\cdots,r\})
\]
这个补集中有$s-r$个数。现在我们假设$s>r$，故此时我们不妨设他们为$n_1,\cdots,n_{s-r}$。于是我们有
\[
a=p_1\cdots p_r=\left(\prod_{i=1}^r q_{\sigma(i)}\right)\cdot\left(\prod_{j=1}^{s-r} q_{n_j}\right)
\]
由整区中的消去律知：
\[
1\sim \prod_{j=1}^{s-r} q_{n_j}
\]
于是$\forall 1\leq j\leq s-r$都有$q_{n_j}$是可逆元，即单位。这与$q_{n_j}$是不可约元矛盾，因此必有$s=r$。\\
\indent 此时我们可知$\sigma$是一个双射。而当$s=r$时，特别地有$\{1,\cdots,s\}=\{1,\cdots,r\}$是同一个集合，因此$\sigma$是$\{1,\cdots,r\}$上的置换，即$\sigma\in S_r$，且有：
\[
\forall 1\leq i\leq r,\; p_i\sim q_{\sigma(i)}
\]
因此可知$R$是唯一析因环。
\end{proof}

\section{主理想整环PID和欧几里得整环ED}

\indent 回忆之前定义过，可以由单个元素生成的理想我们称之为主理想。\\
\kfj{定义：}设$R$是整区，若$R$中所有的理想都是主理想，即都可由单个理想生成，则称$R$为主理想整环(Principal Ideal Domain)，简记为PID。\\
\noteb{(注：由于PID是整区，因此PID中的理想都形如：$<a>=\{ra|r\in R\}$，且$x\in <a>\iff a|x$)}\\
\kfj{例：}$\mathbb{Z}$是PID。这是因为$\mathbb{Z}$的子群都形如$m\mathbb{Z}=<m>$从而必然是主理想。\\
\kfj{例：}整系数多项式$\mathbb{Z}[x]$在多项式加法和多项式乘法下构成一个整区，但是$\mathbb{Z}[x]$不是PID。
\begin{proof}
显然$\mathbb{Z}[x]$是一个整区，下面说明它不是PID。\\
\indent 考虑由$2$和$x^2+1$生成的理想$I=<2,x^2+1>$，假设$I$是$\mathbb{Z}[x]$的主理想，则存在多项式$g(x)\in \mathbb{Z}[x]$，使得$I=<g(x)>$。于是可知$2\in I\implies g(x)\mid 2$，易知$2$是$\mathbb{Z}[x]$中的不可约元素，从而可知$g(x)\in \{1,-1,2,-2\}$。然而回忆含幺交换环的生成理想，可知：
\[
<2,x^2+1>=\{2u(x)+(x^2+1)v(x)\mid u(x),v(x)\in \mathbb{Z}[x]\}
\]
\indent 因此$g(x)\neq 1$，同理$g(x)\neq -1$，故$g(x)\in \{2,-2\}$。 显然$x^2+1$也是$\mathbb{Z}[x]$中的不可约元，从而与$g(x)\in \{2,-2\}$矛盾。综上可知$I$不是主理想，因而$\mathbb{Z}[x]$不是PID。
\end{proof}
\kfj{定理：}PID必定是UFD，即主理想整环一定是唯一析因环。
\begin{proof}
回忆$UFD\iff $(因子链条件)$+$(最大公因子条件)，我们先证明PID满足因子链条件。\\
\indent 设$R$是PID，$a_1,a_2,a_3,\cdots$是$R$中的一条因子链，则由整除关系知：
\[
<a_1>\subset<a_2>\subset<a_3>\subset\cdots
\]
记：
\[
I=\bigcup_{k\geq 1}<a_k>
\]
\indent 则$\forall x\in I,\exists k\geq 1,x\in <a_k>\implies \forall r\in R, rx\in <a_k>\subset I$，从而可知$I$是$R$的理想。由于$R$是PID，因此不妨设$I=<a>$\\
\indent 由于$a\in I$，故存在$m\geq 1$，使得$a\in <a_m>$。因此可知$I=<a_m>$，于是：
\[
\forall n\geq m,\; <a_m>\subset<a_n>\subset I\implies <a_m>=<a_n>=I\iff a_n\sim a_m
\]
因此可知$R$满足因子链条件。下面说明$R$满足最大公因子条件。\\
\indent 任取$a,b\in R^*$，作$a$和$b$的生成理想$<a,b>\zg R$。由于$R$是PID，故$\exists d\in R^*$，使得$<d>=<a,b>$。此时由$a\in <d>$，且$b\in <d>$可知：$d\mid a,d\mid b$，即$d$是$a$和$b$的公因子。若$d'$也是$a$和$b$的公因子，则必有$d'\mid a$且$d'\mid b$，从而由生成理想的极小性知：
\[
<d>=<a,b>\subset<d'>\implies d\mid d'
\]
因此可知$d$是$a$和$b$的最大公因子，从而$R$满足最大公因子条件。综上可知，$R$是UFD。
\end{proof}
\kfj{推论：}(Bezout等式)在PID中，任意两个元素$a$和$b$，考虑$d$是$a$和$b$的最大公因子，则：
\[
\exists u,v\in R,\; \text{使得}ua+vb=d
\]
\begin{proof}
已知$<d>=<a,b>=\{xa+yb|x,y\in R\}$，于是由$d\in <a,b>$立即得出结论。
\end{proof}
\noteb{(注：虽然我们在上述定理中只证明了$<a,b>$可由二者的最大公因子生成，但由于同一个理想的生成元是互相相伴的，且最大公因子之间也是互相相伴的，因此$<a,b>$的生成元与二者的所有最大公因子一一对应)}\\
\indent 我们不难看出，整数环是整区这个对象的一个代表，一般的整区与整数环有诸多相似之处，上面我们就推广了某些整区上可以成立Bezout等式，这是初等数论中一个比较广为人知的结论。事实上，一般的整区与整数环最大的相似之处就是没有普遍成立的除法。但是初等数论给出整数环中除法的替代品：带余除法。事实上这一特殊的操作并不能在一般的整区上推广，我们把这个性质提炼出来，称这种”能做带余除法“的整区称为欧几里得整环，即欧式整环。\\
\kfj{定义：}对于整区$R$，记$\mathbb{N}$是自然数集\noteb{($\{0,1,2,\cdots\}$)}，若存在映射
\[
\delta\colon R^*\to \mathbb{N}
\]
使得$\forall (a,b)\in R\times R^*$，满足：
\[\left\{
\begin{aligned}
&(1)\text{要么}\exists q\in R,\; a=qb\\
&(2)\text{要么}\exists q,r\in R,\;a=qb+r,\; r\neq 0,\; \delta(r)<\delta(b)\\
\end{aligned}\right.\]
则称$R$是一个欧几里得整环，简记为ED(Euclidean Domain)\\
\kfj{例：}$\mathbb{Z}$是ED。此处令$\delta\colon \mathbb{Z}-\{0\}\to \mathbb{N},\; n\mapsto |n|$，则易知带余除法成立。事实上，如此定义的带余除法就算是在整数环上也不见得唯一。
\[
\left\{
\begin{aligned}
&5=2\times 2+1,\; \delta(1)<\delta(2)\\
&5=3\times 2+(-1),\; \delta(-1)<\delta(2)
\end{aligned}\right.
\]
下面给出一个比较重要的例子：\\
\kfj{定理：}设$F$是域，那么$F-$系数多项式环$F[x]$是ED。且对于
\[
\delta\colon F[x]-\{0\}\to \mathbb{N},\; f(x)\mapsto deg(f(x))
\]
有带余除法唯一。
\begin{proof}
(存在性）：\\
\indent 任取两个多项式$f(x),g(x)\in F[x]$，$g(x)$不是零。若$deg(f(x))<deg(g(x))$，则立即有带余除法：
\[
f(x)=0\cdot g(x)+f(x),\; deg(f(x),\; deg(f(x))<deg(g(x))
\]
\indent 若$deg(f(x))\geq deg(g(x))$，我们对多项式的次数使用数学归纳法：设$deg(f(x))=n$。当$n=0$时，可知$f(x)$和$g(x)$都是常多项式，从而由$F$是域知，$g(x)\mid f(x)$，因此也有带余除法。现在对任意$n>0$，归纳假设对所有次数小于$n$的多项式均成立带余除法，则此时设
\[
f(x)=\sum_{k=0}^{n}a_kx^k,\; g(x)=\sum_{s=0}^{m}b_sx^s,\;\forall 0\leq k\leq n,\;\forall 0\leq s\leq m,\; a_k,b_s\in F
\]
则由于$m\leq n$，故$x^{n-m}\in F[x]$，此时做：
\[
f(x)-a_n\inv{b_m}x^{n-m}g(x)=\sum_{k=0}^{n}a_kx^k-\sum_{s=0}^{m}a_n\inv{b_m}b_sx^{n-m+s}
\]
则$f(x)-a_n\inv{b_m}x^{n-m}g(x)$是一个低于$n$次的多项式，于是有带余除法：
\[
f(x)-a_n\inv{b_m}x^{n-m}g(x)=q'(x)g(x)+r(x)
\]
其中$r(x)=0$或$r(x)\neq 0$且$deg(r(x))<deg(g(x))$。于是有带余除法：
\[
f(x)=(q'(x)+a_n\inv{b_m}x^{n-m})g(x)+r(x)
\]
则仍有$r(x)=0$或$r(x)\neq 0$且$deg(r(x))<deg(g(x))$。因此可知带余除法在$F[x]$中总是存在的。\\
(唯一性)：\\
\indent 设$f(x)=q(x)g(x)+r(x)=q'(x)g(x)+r'(x)$，则有：$(q(x)-q'(x))g(x)=r'(x)-r(x)\implies deg((q(x)-q'(x))g(x))=deg(r'(x)-r(x))$。而若$deg(q(x)-q'(x))\geq 0$，则有：$deg((q(x)-q'(x))g(x))\geq deg(g(x))>deg(r'(x)-r(x))$与等式矛盾，因此可知必有$q(x)-q'(x)$是零，从而有$q(x)-q'(x)=r'(x)-r(x)=0$，因此可知带余除法是唯一的。
\end{proof}
\kfj{例：}$\mathbb{Z}[\sqrt{-1}]=\{a+b\sqrt{-1}\mid a,b\in \mathbb{Z}\}$是ED。其中，我们令
\[
\delta\colon \mathbb{Z}[\sqrt{-1}]-\{0\}\to \mathbb{N},\; a+b\sqrt{-1}\mapsto a^2+b^2
\]
下面我们证明此时带余除法成立。
\begin{proof}
任取$\alpha,\beta\in \mathbb{Z}[\sqrt{-1}],\beta\neq 0$，考虑集族:
\[
W(m,n)=\left\{x+y\sqrt{-1}\bigm|-\frac{1}{2}+m\leq x\leq\frac{1}{2}+m,\; -\frac{1}{2}+n\leq y\leq \frac{1}{2}+n\right\}, m,n\in \mathbb{Z}
\]
显然$\{W(m,n)\}_{m,n\in\mathbb{Z}}$是整个复平面$\mathbb{C}$的一个覆盖，因此
\[
\exists q\in \mathbb{Z}[\sqrt{-1}],\; \dfrac{\alpha}{\beta}\in W(Re(q),Im(q))
\]
相应地，有：
\[
|Re(\frac{\alpha}{\beta}-q)|\leq \frac{1}{2},\; |Im(\frac{\alpha}{\beta}-q)|\leq\frac{1}{2}
\]
于是若令$r=\alpha-\beta q$，则有：
\[
\delta(r)={||\alpha-\beta q||}^2={||\beta||}^2{||\frac{\alpha}{\beta}-q||}^2\leq{||\beta||}^2\cdot(\frac{1}{4}+\frac{1}{4})<\delta(\beta)
\]
因此可知当$r=0$时，有整除：$\alpha=\beta q$，否则有带余除法：
\[
\alpha=\beta q+r,\; r\neq 0,\; \delta(r)<\delta(\beta)
\]
故此时$\mathbb{Z}[\sqrt{-1}]$上带余除法成立。
\end{proof}
\kfj{定理：}ED一定是PID
\begin{proof}
设$R$是ED，$I\zg R$是一个理想，下面证明$I$是主理想。\\
若$I$是零理想，则显然$I$是主理想，故不妨设$I$不是零理想。那么考虑非负整数集
\[
S=\{\delta(x)|x\in I-\{0\}\}\subset \mathbb{N}
\]
\indent 那么可知$S$有最小值，于是不妨设$b\in I=\{0\}$，满足$\delta(b)=min(S)$。显然有$<b>\subset I$，下面证明反方向的包含。任取$a\in I$，由于$R$是ED，故做带余除法：
\[
\left\{
\begin{aligned}
& \text{若}\exists q\in R,\; a=qb,\text{则}b|a\implies a\in <b>\\
& \text{若}\exists q,r\in R,\; a=qb+r,\; r\neq 0,\delta(r)<\delta(b),\text{则与}\delta(b)\text{最小矛盾}
\end{aligned}\right.
\]
因此可知必然有：$I=<b>\implies I$是主理想。综上可知，ED一定是PID。
\end{proof}
\indent 我们之前知道，在PID中任意两个元素是存在最大公因子的，但是我们没有一种”算法“找到这个最大公因子。那么下面我们就引入数论中很有名的算法——辗转相除法。\\
\kfj{定理：}ED上可使用辗转相除法求最大公因子。
\begin{proof}
\indent 我们先来说明可以通过这个算法来找到公因子。$\forall a,b\in R^*$，不妨设$\delta(a)\geq \delta(b)$，若$b\mid a$，则可知$b$是$a$和$b$的公因子。否则由带余除法：
\[
\exists q_1,r_1\in R,\text{使得}a=q_1b+r_1,\; r_1\neq 0,\; \delta(r_1)<\delta(b)
\]
同理若$r_1\mid b$，那么可知$r_1$是$a$和$b$的公因子，否则由带余除法：
\[
\exists q_2,r_2\in R,\text{使得}b=q_2r_1+r_2,\; r_2\neq 0,\; \delta(r_2)<\delta(r_1)
\]
这个操作可以一直重复下去，得到：（假设一直不可整除）
\[
\left\{
\begin{aligned}
&a=q_1b+r_1,\; r_1\neq 0 ,\; \delta(r_1)<\delta(b)\\
&b=q_2r_1+r_2,\; r_2\neq 0,\; \delta(r_2)<\delta(r_1)\\
&r_1=q_3r_2+r_3,\; r_3\neq 0,\; \delta(r_3)<\delta(r_2)\\
&\vdots\\
&r_k=q_{k+2}r_{k+1}+r_{k+2},\; r_{k+2}\neq 0,\; \delta(r_{k+2})<\delta(r_{k+1})
\end{aligned}\right.
\]
若此操作无法在有限步内停止，则得到无穷序列：
\[
\delta(b)>\delta(r_1)>\delta(r_2)>\delta(r_3)>\cdots
\]
这与$Im\delta\subset\mathbb{N}$必有下界矛盾，因此存在正整数$n$，使得$r_{n-1}=q_{n+1}r_n$\\
\indent 可知$r_n|r_{n-1}$，因此$r_{n-2}=q_nr_{n-1}+r_n\implies r_n|r_{n-2}$，以此类推，由$r_k=q_{k+2}r_{k+1}+r_{k+2}$知：$r_n\mid r_1$。同理由$b=q_2r_1+r_2,a=q_1b+r_1$知$r_n\mid a$，且$r_n\mid b$。因此我们如此求出$r_n$是$a$和$b$的公因子，下面证明$r_n$也是它们的最大公因子。\\
\indent 若$d\mid a,\; d\mid b$，则由$r_1=a-q_1b$知$d\mid r_1$。同理再由$r_2=b-q_2r_1,\; r_{k+2}=r_k-q_{k+2}r_{k+1}$知，$d\mid r_n$，因此$r_n$是$a$和$b$的最大公因子。
\end{proof}

\section{素理想与极大理想：特殊的理想给出特殊的商环}
\kfj{定义：}$R$是环，$\mathfrak{m}\subset R$是真理想。若包含$\mathfrak{m}$的理想仅有$\mathfrak{m}$和$R$，则称$\mathfrak{m}$是一个极大理想（Maximal Ideal）。\\
\indent \kfj{引理：}对于含幺环$R$，有：$R$是除环$\iff\{0\}$是$R$中的极大理想。
\begin{proof}
$(\implies)$若$R$是除环，则设理想$M$是非零理想，那么$\exists m\in M,\; m\neq 0$，于是可知$\forall r\in R^*,\; r=r\inv{m}\cdot m\in <m>$，因此可知$M=R$。从而有$\{0\}$是极大理想。\\
$(\impliedby)$若$\{0\}$是极大理想，那么$\forall r\in R^*$，可知$\{0\}\subsetneq <r>$。由$\{0\}$是极大理想知$<r>=R$，特别地$r$是可逆元。于是由$r$的任意性可知$R$是除环。
\end{proof}
\kfj{定理：}$R$是含幺环，对于理想$\mathfrak{m}\subsetneq R$，有：
\[
\mathfrak{m}\text{是}R\text{的极大理想}\iff R/\mathfrak{m}\text{是除环}
\]
\begin{proof}
由引理已知$R/\mathfrak{m}$是除环$\iff\{\ovl{0}\}$是$R/\mathfrak{m}$的极大理想$\iff$包含$\{\ovl{0}\}$的理想只有$\{\ovl{0}\}$和$R/\mathfrak{m}$。对于自然同态$\pi\colon R\to R/\mathfrak{m}$，考虑满同态对应定理，可知包含$\mathfrak{m}$的理想只有$\mathfrak{m}$和$R$。因此可知$R/\mathfrak{m}$是除环$\iff\mathfrak{m}$是$R$的极大理想。
\end{proof}
从这个定理我们可以总结出一个证明理想的极大性的一个套路，也就是把这个理想实现为打到除环的满同态的核。\\
\kfj{例：}$<n>$是$\mathbb{Z}$的极大理想$\iff n$是素数。
\begin{proof}
可知$\mathbb{Z}/<m>\cong \mathbb{Z}_n$，故由定理可知$<m>$是$\mathbb{Z}$的极大理想$\iff\mathbb{Z}_n$是域。而$\mathbb{Z}_n$是域等价于$n$是素数。因此$<m>$是$\mathbb{Z}$的极大理想等价于$m$是素数。
\end{proof}
\indent \kfj{例：}$a\in \mathbb{C}$，则在$\mathbb{C}[x]$中\[<x-a>=Ker(\phi\colon \mathbb{C}[x]\to \mathbb{C},\; f(x)\mapsto f(a) )\]
从而是极大理想。\\
\indent \kfj{例：}$n\in \mathbb{Z}$，$p$是素数，则在$\mathbb{Z}[x]$中
\[
<p,(x-n)>=Ker(\phi\colon \mathbb{Z}[x]\to \mathbb{Z}_p,\; f(x)\mapsto f(\ovl{n}) )
\]
从而是极大理想。

\kfj{定义：}$R$是一个环，若真理想$\mathfrak{p}\subsetneq R$满足：
\[
ab\in \mathfrak{p}\implies a\in \mathfrak{p}\text{或}b\in \mathfrak{p}
\]
则称$\mathfrak{p}$是一个素理想。\\
\indent 不难联想到，素理想的定义跟素元素的定义极其相似，我们将会看到这二者之间的联系。\\
\noteb{(注：我们考虑上述定义的逆否命题，$\mathfrak{p}$是素理想$\iff a\notin \mathfrak{p}$且$b\notin\mathfrak{p}\implies ab\notin\mathfrak{p}$，即集合$R-\mathfrak{p}$是乘法封闭的。当然这个集合只能确定的是乘法封闭，不会继承$R$更多的结构。)}\\
\kfj{定理：}环$R$中真理想$\mathfrak{p}$是素理想$\iff R/\mathfrak{p}$是整的。
\begin{proof}
对于自然同态$\pi\colon R\to R/\mathfrak{p}$，可知$x\notin \mathfrak{p}\iff \pi(x)\neq \ovl{0}$于是有：
\[
ab\notin \mathfrak{p}\iff \pi(ab)\neq \ovl{0}\iff \pi(a)\pi(b)\neq \ovl{0}
\]
于是可知$\mathfrak{p}$是素理想等价于任意$a,b\notin \mathfrak{p}$，都满足$ab\notin \mathfrak{p}$，即$\pi(a)\pi(b)\neq \ovl{0}$。再由$\pi$是满同态就可以得到$\mathfrak{p}$是素理想等价于$R/\mathfrak{p}$是整的。
\end{proof}
\kfj{推论：}含幺环中，极大理想一定是素理想。
\begin{proof}
我们考虑$\mathfrak{m}$是环$R$的极大理想，则可知$R/\mathfrak{m}$是域，特别地是整的，因此由定理知$\mathfrak{m}$是素理想。
\end{proof}
\noteb{(注：如果去掉“含幺”这个条件，则命题不成立。)}\\
\kfj{例：}考虑环$2\mathbb{Z}$的理想$4\mathbb{Z}$，那么可知$2\mathbb{Z}$和$4\mathbb{Z}$之间没有任何其他理想，从而$4\mathbb{Z}$是$2\mathbb{Z}$的极大理想，但是由于
\[
2\notin 4\mathbb{Z},\; 2\cdot 2=4\in 4\mathbb{Z}
\]
因此$4\mathbb{Z}$不是$2\mathbb{Z}$中的素理想。\\
\kfj{例：}$<n>$是$\mathbb{Z}$的素理想$\iff n$是素数。
\begin{proof}
可知$\mathbb{Z}/<m>\cong \mathbb{Z}_n$，故由定理可知$<m>$是$\mathbb{Z}$的极大理想$\iff\mathbb{Z}_n$是整的。$n$是素数时显然有$\mathbb{Z}_n$是整的，但$n$不是素数时其真因子就是$\mathbb{Z}_n$中的零因子，从而不是整的。因此$<m>$是$\mathbb{Z}$的极大理想等价于$m$是素数。
\end{proof}
这个例子说理想的“素性”与元素的“素性”有强联系（事实上这种相似性在定义上就已经体现出来了），那么我们是否可以把这个命题推广到一般的环上呢？下面的定理给出了回答。\\
\kfj{定理：}对于整区$R$，有：$a$是素元素$\iff<a>$是素理想
\begin{proof}
\[
<a>\text{是素理想}\iff \left\{
\begin{aligned}
&<a>\neq R\\
&x,y\notin<a>\implies xy\notin<a>
\end{aligned}\right.
\iff\left\{
\begin{aligned}
&a\in R-U\\
&a\nmid x,\; a\nmid y\implies a\nmid xy
\end{aligned}\right.
\iff a\text{是素元素}
\]
\end{proof}
\kfj{推论：}由于UFD满足(素性条件)，因此由不可约元等价于素元素知：
\[
<a>\text{是素理想}\iff a\text{是不可约元}
\]
\kfj{定理：}若$R$是PID，$I$是$R$的某个理想，则$I$是极大理想$\iff I$是素理想。
\begin{proof}
由于$R$是PID，那么$I$是主理想，于是不妨设$I=<x>$。于是考虑另一个理想$A$，同理记$A=<a>$，那么有：\[
I\subset A\subset R\iff<x>\subset <a>\iff a\mid x
\]
因此可知$I$是$R$的极大理想等价于当$a\mid x$时，必有$a\in U$或$a\sim x$。故等价于$x$不可约。由于PID一定是UFD，因此可知等价于$x$是素元。于是可知$I$是极大理想等价于$I$是素理想。
\end{proof}
下面我们给出一种常见的，构造域的方式。\\
\kfj{例：}对于域系数多项式环$F[x]$，考虑其中的不可约多项式$f(x)$。可知$F[x]$是ED从而是PID。因此由上述定理可知$<f(x)>$是极大理想，因此$F[x]/<f(x)>$是域。\\
\kfj{例：}构造125阶域。（这个例子需要一些许线性代数基础，学过线性代数的读者可以看一下）。\\
\indent 考虑$\mathbb{Z}_5$是域，那么$\mathbb{Z}_5[x]$是PID。考虑多项式：
\[
f(x)=x^3+x^2+1\in \mathbb{Z}_5[x]
\]
可知$degf(x)=3$，且$f(0)=1,f(1)=3,f(2)=3,f(3)=2,f(4)=1$，于是$f$是不可约三次多项式。故有
\[
K=\mathbb{Z}_5[x]/<x^3+x^2+1>\text{是域}
\]
由带余除法知，商环K中的等价类一定可以由2次多项式代表。于是有
\[
K=\{a_0+a_1x+a_2x^2|a_0,a_1,a_2\in \mathbb{Z}_5\}
\]
\indent 显然有$\{1,x,x^2\}$在K中线性无关，因此可知$K$是$\mathbb{Z}_5$上的三维线性空间。我们写出$K$在$\{1,x,x^2\}$下的基的坐标表示：$K=\{(a_0,a_1,a_2)|a_0,a_1,a_2\in \mathbb{Z}_5\}$。故$K$的元素个数为：$5^3=125$。